\documentclass[12pt,a4paper]{article}

\usepackage[utf8]{inputenc}
\usepackage[T1]{fontenc}
\usepackage[provide=*,spanish]{babel}
\usepackage{mathpazo}
\usepackage{microtype}
\usepackage{amsmath,amssymb,amsthm}
\usepackage{booktabs}
\usepackage{tabularx}
\usepackage{geometry}
\usepackage{enumitem}
\usepackage[hidelinks]{hyperref}

\geometry{margin=2.5cm}

\title{Ecuaciones Diferenciales Ordinarias\\Clase IV -- Teórico}
\author{Benjamín Marcolongo}
\date{}

\begin{document}

	\maketitle

	\section*{Introducción}

	En la Clase III estudiamos integración directa y separación de variables. Sin embargo, no toda ecuación de primer orden puede resolverse separando las variables. En esta clase estudiaremos las ecuaciones lineales de primer orden.

	El factor integrante aparecerá como el método general para resolverlas. Su función es transformar el miembro izquierdo de la ecuación en la derivada de un producto, que luego puede integrarse directamente. Esta construcción también permitirá demostrar el teorema de existencia y unicidad incluido en el apéndice.

	En todo el material, \(x\) representa la variable independiente e \(y=y(x)\) la función desconocida. Escribiremos \(dy/dx\) para la derivada de \(y\) respecto de \(x\).

	\clearpage
	\section{Ecuaciones lineales de primer orden}

	Una ecuación diferencial de primer orden es \textbf{lineal en \(y\)} si puede escribirse como

	\[
	a_1(x)\frac{dy}{dx}+a_0(x)y=g(x),
	\]
	%
	donde \(a_1\), \(a_0\) y \(g\) son funciones conocidas de \(x\). En un intervalo \(I\) donde \(a_1(x)\neq0\), podemos dividir por \(a_1(x)\) y obtener la \textbf{forma estándar}

	\[
	\boxed{
	\frac{dy}{dx}+p(x)y=q(x),
	}
	\qquad
	p(x)=\frac{a_0(x)}{a_1(x)},
	\qquad
	q(x)=\frac{g(x)}{a_1(x)}.
	\]
	%
	La ecuación es homogénea si \(q(x)=0\) en \(I\), y es no homogénea si \(q\) no es idénticamente nula.

	Los ceros de \(a_1\) son puntos singulares de la ecuación: la forma estándar no está definida allí. Por eso, el método debe aplicarse en un intervalo que no atraviese esos puntos. Si \(p\) y \(q\) son continuas en \(I\), entonces para cada condición inicial

	\[
	y(x_0)=y_0,
	\qquad x_0\in I,
	\]
	%
	existe una única solución definida en todo el intervalo \(I\). La demostración de este resultado y algunos comentarios sobre sus hipótesis se presentan en el Apéndice~\ref{ap:existencia-unicidad-lineales}.

	\subsection{Construcción del factor integrante}

	Consideremos la ecuación lineal en forma estándar

	\[
	\frac{dy}{dx}+p(x)y=q(x).
	\]
	%
	Buscamos una función no nula \(\mu=\mu(x)\) tal que, al multiplicar toda la ecuación por ella, el miembro izquierdo se convierta en la derivada de un producto:

	\[
	\mu(x)\frac{dy}{dx}+\mu(x)p(x)y
	=
	\frac{d}{dx}\bigl(\mu(x)y(x)\bigr).
	\]
	%
	Por la regla del producto,

	\[
	\frac{d}{dx}\bigl(\mu y\bigr)
	=
	\mu\frac{dy}{dx}+\frac{d\mu}{dx}y.
	\]
	%
	Por lo tanto, necesitamos que

	\[
	\frac{d\mu}{dx}=p(x)\mu(x).
	\]
	%
	Esta ecuación es separable. En cualquier intervalo donde \(p\) sea continua, podemos elegir

	\[
	\boxed{
	\mu(x)=\exp\left(\int p(x)\,dx\right).
	}
	\]
	%
	La función \(\mu\) se denomina \textbf{factor integrante}. No es necesario agregar una constante al calcular la integral: multiplicar \(\mu\) por una constante no nula no modifica el método.

	Al multiplicar la ecuación por \(\mu\), obtenemos

	\[
	\frac{d}{dx}\bigl(\mu(x)y(x)\bigr)
	=
	\mu(x)q(x).
	\]
	%
	Integrando respecto de \(x\),

	\[
	\mu(x)y(x)
	=
	\int \mu(x)q(x)\,dx+C,
	\]
	%
	y finalmente

	\[
	\boxed{
	y(x)
	=
	\frac{1}{\mu(x)}
	\left[
	\int \mu(x)q(x)\,dx+C
	\right].
	}
	\]

	\begin{quote}
		\textbf{Importante.} El factor integrante se calcula después de escribir la ecuación en forma estándar. Si se omite la división por \(a_1(x)\), se obtiene en general un factor incorrecto.
	\end{quote}

	\subsection{Procedimiento}

	Para resolver una ecuación lineal de primer orden conviene seguir este orden:

	\begin{enumerate}[label=\arabic*.]
		\item identificar un intervalo \(I\) donde el coeficiente principal no se anule;
		\item escribir la ecuación en la forma estándar;
		\item identificar \(p(x)\) y \(q(x)\);
		\item calcular el factor integrante \(\mu(x)\);
		\item multiplicar toda la ecuación por \(\mu(x)\);
		\item reconocer la derivada del producto \(\mu(x)y(x)\);
		\item integrar e imponer la condición inicial, si existe;
		\item determinar el intervalo de definición de la solución.
	\end{enumerate}

	\clearpage
	\subsection{Ejemplo: el coeficiente principal se anula}

	Consideremos

	\[
	2x\frac{dy}{dx}-3y=2x^2.
	\]
	%
	El coeficiente de \(dy/dx\) se anula en \(x=0\). Trabajaremos primero en el intervalo \(x>0\). Al dividir por \(2x\),

	\[
	\frac{dy}{dx}-\frac{3}{2x}y=x.
	\]
	%
	En este intervalo,

	\[
	p(x)=-\frac{3}{2x},
	\qquad
	q(x)=x,
	\]
	%
	y el factor integrante es

	\[
	\mu(x)
	=
	\exp\left(-\frac32\int\frac{1}{x}\,dx\right)
	=
	x^{-3/2}.
	\]
	%
	Multiplicando la ecuación por \(x^{-3/2}\),

	\[
	x^{-3/2}\frac{dy}{dx}
	-
	\frac32x^{-5/2}y
	=
	x^{-1/2},
	\]
	%
	por lo que

	\[
	\frac{d}{dx}\left(x^{-3/2}y\right)=x^{-1/2}.
	\]
	%
	Integrando,

	\[
	x^{-3/2}y=2x^{1/2}+C,
	\]
	%
	y entonces

	\[
	\boxed{y(x)=2x^2+Cx^{3/2},\qquad x>0.}
	\]
	%
	En un intervalo contenido en \(x<0\) debe repetirse el análisis usando \(\ln|x|\). Una solución no puede describirse automáticamente mediante una única constante a través del punto singular \(x=0\).

	\subsection{Ejemplo: una integral no elemental}

	Consideremos

	\[
	\frac{dy}{dx}+2xy=1.
	\]
	%
	Aquí \(p(x)=2x\) y \(q(x)=1\), ambas continuas en toda la recta real. El factor integrante es

	\[
	\mu(x)=\exp\left(\int 2x\,dx\right)=e^{x^2}.
	\]
	%
	Por lo tanto,

	\[
	\frac{d}{dx}\left(e^{x^2}y\right)=e^{x^2}.
	\]
	%
	La integral de \(e^{x^2}\) no puede expresarse mediante funciones elementales. Esto no impide escribir la solución general. Fijando un punto de referencia \(x_0\in\mathbb{R}\),

	\[
	\boxed{
	y(x)
	=
	e^{-x^2}
	\left[
	C+\int_{x_0}^{x}e^{s^2}\,ds
	\right],
	}
	\]
	%
	donde \(s\) es la variable muda de integración y \(C\) es una constante real.

	\section{Cómo reconocer el método}

	\begin{center}
	\begin{tabularx}{\textwidth}{@{}
		>{\raggedright\arraybackslash}p{0.22\textwidth}
		>{\raggedright\arraybackslash}X
		>{\raggedright\arraybackslash}X@{}}
		\toprule
		\textbf{Tipo} & \textbf{Forma característica} & \textbf{Idea principal} \\
		\midrule
		Lineal &
		\(\dfrac{dy}{dx}+p(x)y=q(x)\) &
		Multiplicar por \(\mu(x)=\exp(\int p(x)\,dx)\). \\
		\bottomrule
	\end{tabularx}
	\end{center}

	Antes de aplicar el factor integrante conviene revisar si la ecuación también es separable, pues en algunos casos ese método puede ser más simple. Siempre deben controlarse las divisiones realizadas, los puntos singulares y el intervalo donde la solución es válida.

	\section*{Articulación con la Guía 1}

	El método desarrollado permite resolver el problema 21 de la Guía 1.

	\section*{Bibliografía y recursos recomendados}

	\begin{itemize}[leftmargin=*]
		\item D. G. Zill, \textit{A First Course in Differential Equations with Modeling Applications}, 12.\textsuperscript{a} ed., sección 2.3.
	\end{itemize}

	\clearpage
	\appendix
	\section{Existencia y unicidad}
	\label{ap:existencia-unicidad-lineales}

	\subsection{Enunciado}

	Sea \(I\) un intervalo abierto de la recta real y sean \(p:I\to\mathbb{R}\) y \(q:I\to\mathbb{R}\) funciones continuas. Dados un punto \(x_0\in I\) y un valor \(y_0\in\mathbb{R}\), el problema de valor inicial

	\[
	\frac{dy}{dx}+p(x)y(x)=q(x),
	\qquad
	y(x_0)=y_0,
	\]
	%
	posee una única solución \(y:I\to\mathbb{R}\). En particular, la solución está definida en todo el intervalo \(I\), no solamente en un entorno del punto inicial \(x_0\).

	\subsection{Demostración de la existencia}

	Definimos el factor integrante normalizado en el punto inicial mediante

	\[
	\mu(x)
	=
	\exp\left(\int_{x_0}^{x}p(s)\,ds\right),
	\]
	%
	donde \(s\) es una variable muda de integración. Como \(p\) es continua en \(I\), la integral existe para todo \(x\in I\). Además, \(\mu(x)>0\) y

	\[
	\frac{d\mu}{dx}=p(x)\mu(x),
	\qquad
	\mu(x_0)=1.
	\]
	%
	Multiplicando la ecuación diferencial por \(\mu(x)\) y usando la regla del producto, obtenemos

	\[
	\frac{d}{dx}\bigl(\mu(x)y(x)\bigr)
	=
	\mu(x)q(x).
	\]
	%
	Integramos entre \(x_0\) y \(x\):

	\[
	\mu(x)y(x)-\mu(x_0)y(x_0)
	=
	\int_{x_0}^{x}\mu(s)q(s)\,ds.
	\]
	%
	Al utilizar \(\mu(x_0)=1\) e \(y(x_0)=y_0\), resulta

	\[
	\boxed{
	y(x)
	=
	\frac{1}{\mu(x)}
	\left[
	 y_0+\int_{x_0}^{x}\mu(s)q(s)\,ds
	\right].
	}
	\]
	%
	Esta expresión está bien definida para todo \(x\in I\), porque \(\mu(x)\neq0\) y el producto \(\mu q\) es continuo. Al derivarla se recuperan la ecuación diferencial y la condición inicial. Por lo tanto, la fórmula construye una solución en todo \(I\).

	\subsection{Demostración de la unicidad}

	Supongamos que \(y_1:I\to\mathbb{R}\) e \(y_2:I\to\mathbb{R}\) son dos soluciones del mismo problema de valor inicial. Definimos

	\[
	z(x)=y_1(x)-y_2(x).
	\]
	%
	Al restar las ecuaciones satisfechas por \(y_1\) e \(y_2\), obtenemos

	\[
	\frac{dz}{dx}+p(x)z(x)=0,
	\qquad
	z(x_0)=0.
	\]
	%
	Multiplicando por el mismo factor integrante \(\mu(x)\), resulta

	\[
	\frac{d}{dx}\bigl(\mu(x)z(x)\bigr)=0.
	\]
	%
	Por lo tanto, \(\mu(x)z(x)\) es constante en \(I\). Como \(\mu(x_0)z(x_0)=0\), esa constante es cero. Dado que \(\mu(x)>0\), se sigue que

	\[
	z(x)=0
	\]
	%
	para todo \(x\in I\). En consecuencia, \(y_1(x)=y_2(x)\) en todo \(I\), lo que demuestra la unicidad.

	\subsection{Comentarios sobre las hipótesis}

	\begin{itemize}[leftmargin=*]
		\item La continuidad de \(p\) garantiza que el factor integrante esté definido y sea derivable en todo \(I\). La continuidad de \(q\) garantiza que \(\mu q\) pueda integrarse y que la función obtenida sea una solución clásica.
		\item Si la ecuación original es \(a_1(x)y'(x)+a_0(x)y(x)=g(x)\), el intervalo \(I\) debe elegirse de modo que \(a_1(x)\neq0\) en todo \(I\). Además, los cocientes \(p=a_0/a_1\) y \(q=g/a_1\) deben ser continuos allí. El teorema no permite atravesar automáticamente un cero de \(a_1\).
		\item El resultado también puede relacionarse con el teorema de Picard--Lindelöf. Al escribir la ecuación como \(dy/dx=q(x)-p(x)y\), el miembro derecho es continuo en \(x\) y localmente Lipschitz respecto de \(y\). Picard--Lindelöf proporciona existencia y unicidad local; la fórmula explícita obtenida con el factor integrante muestra, además, que la solución se extiende a todo el intervalo \(I\).
		\item Para demostrar este teorema lineal no es necesario invocar un resultado general de existencia y unicidad: el propio método del factor integrante construye la solución y permite probar directamente que no puede haber otra.
	\end{itemize}

	\subsection*{Nota sobre la elaboración del material}

	Este material fue elaborado por Benjamín Marcolongo con la asistencia de \textbf{GPT-5.6 Sol}, un modelo de OpenAI utilizado a través del entorno Codex, para la organización, redacción y revisión del contenido.

\end{document}
